\section{Introduction}

Computational analysis of Greek parliamentary archives from the democratic transition period is constrained by three coupled bottlenecks: noisy OCR, register mismatch between Katharevousa and modern NLP resources, and limited gold-standard syntactic data \cite{barmpounis2024,fitsilis2024}. While prior work has demonstrated that hybrid processing can recover useful stylometric and discourse-level signals, a dedicated and reproducible dependency parsing pipeline for Katharevousa has remained missing.

This study addresses that gap by moving from exploratory hybrid analysis to full parser development with explicit benchmark governance. We frame the contribution not as a single-model result, but as a pipeline outcome:
\begin{itemize}[leftmargin=1.25em]
  \item robust data reconstruction and annotation operations,
  \item controlled train/test split management,
  \item multi-family model comparison under shared metrics,
  \item transparent reporting of failed and successful optimization paths.
\end{itemize}

The goal is twofold: (1) produce a release-ready baseline for open distribution and (2) provide a replicable methodological blueprint for low-resource historical language NLP.
